\section{Atenuación aparente y comparación con la teoría}
A 3 GHz, el flujo radial compensado decayó más rápidamente en agua salada que en tejido muscular y en concreto. Este orden coincidió con el previsto a partir de la permitividad y la conductividad asignadas a los medios (Cuadro~\ref{tab:fdtd-base}). La compensación por $r^2$ permitió examinar el decaimiento que persistió después de descontar la expansión geométrica esperada en la región de radiación. Sin embargo, la pendiente se estimó sobre una línea de observación y en un intervalo finito de distancias. Por ello se interpretó como una atenuación aparente de la radiación del monopolo y se tomó la atenuación teórica del medio homogéneo como referencia de propagación.

En los tres medios con pérdidas, la simulación subestimó la atenuación teórica. El concreto presentó la menor discrepancia relativa (1.31\,\%), mientras que en agua salada y músculo se obtuvieron diferencias de 4.29\,\% y 5.12\,\%, respectivamente. El orden de estas discrepancias difirió del orden de pérdidas de los materiales. Además, la longitud del monopolo y las dimensiones del dominio se escalaron con la longitud de onda de cada medio. La comparación describió, por tanto, configuraciones eléctricamente semejantes, con distinta longitud física y con un diámetro de conductor fijo. Atribuir las diferencias únicamente a la conductividad habría excedido el alcance de este diseño.

\section{Precisión del ajuste y sensibilidad a la ventana radial}
En tejido muscular a 5 px/mm, el IC del 95\,\% abarcó de 0.4663 a 0.4674 dB/mm, mientras que la separación entre la estimación central y la referencia alcanzó 0.0252 dB/mm. La estrechez del intervalo indicó que se recuperó una pendiente semejante al remuestrear bloques de residuos dentro de la ventana elegida. La discrepancia con la teoría incorporó además la influencia del perfil físico, de su discretización y del modelo empleado para extraer la pendiente. Los valores de Durbin--Watson próximos a cero y la curvatura visible de los residuos señalaron que los puntos radiales contiguos compartieron desviaciones del mismo signo. Esta estructura limitó la interpretación de la precisión estadística como medida de exactitud física.

El barrido del límite inferior permitió examinar esa distinción sin repetir la simulación. En músculo a 5 px/mm, al desplazar el inicio de la ventana de $1.5\lambda$ a $3\lambda$, la estimación con el modelo lineal M0 pasó de 0.4650 a 0.4695 dB/mm. La variación de 0.0045 dB/mm superó la anchura del intervalo obtenido en la ventana principal, y aun el mayor valor permaneció por debajo de la referencia de 0.4921 dB/mm. Así, la ventana explicó parte de la variación de la pendiente. El barrido por sí solo dejó abierta la contribución relativa del campo cercano y de las reflexiones de frontera.

\section{Alcance del refinamiento y del control en vacío}
En agua salada, al pasar de 5 a 40 px/mm, la discrepancia respecto de la teoría descendió de 4.29\,\% a 0.81\,\%. El sentido del cambio y la reducción de los incrementos sucesivos indicaron que el refinamiento mejoró la concordancia. Sin embargo, entre 10 y 20 px/mm la selección pasó de M0 a M1, que incorporó un término sinusoidal. Por ello, la diferencia de 3.51\,\% entre las pendientes de las mallas extremas caracterizó la sensibilidad del procedimiento completo. La variación del orden aparente entre tríos de mallas limitó la extrapolación a paso espacial nulo (Cuadro~\ref{tab:fdtd-mallas}). Los perfiles adicionales se reunieron en el \hyperref[an:fdtd-perfiles]{Anexo B}.

A 5 px/mm, incluso la longitud de onda más corta, la del agua salada, abarcó unas 56 celdas. En cambio, el diámetro del monopolo ocupó sólo dos en todos los medios. Esta diferencia permitió distinguir la resolución de la onda propagada de la representación del conductor. La malla común resultó útil para comparar el orden de atenuación y cuantificar cuánto cambió cada pendiente al refinarla. Las diferencias observadas en concreto y músculo, de aproximadamente 0.26\,\% y 3.59\,\% frente a sus respectivas mallas más finas, delimitaron esa comparación para cada material.

En vacío, la ausencia de pérdidas proporcionó un control particularmente exigente. La estimación alcanzó el límite inferior permitido por el ajuste, con $R^2=0.0000$ y un RMSE de 5.8227 dB. El perfil compensado presentó oscilaciones intensas y la coincidencia numérica con la atenuación teórica nula tuvo escaso valor probatorio. Con una sola malla quedó abierta la distinción entre efectos de frontera y de discretización. Por último, el análisis radial cuantificó el decaimiento espacial del flujo observado. Para establecer qué fracción de la potencia se disipó en el medio se requirió una comprobación independiente del balance sobre un volumen de control.

\section{Interpretación del balance local en agua salada}
La tapa inferior aportó potencia al volumen de control (Cuadro~\ref{tab:balance-agua-cinco}). Si sólo se hubiera contado la entrada radial, la disipación habría representado el 100.27\,\% de esa entrada. Al incorporar el transporte axial a través de la tapa inferior, la fracción disipada fue del 81.17\,\% de la entrada total. Así se explicó, mediante las direcciones de flujo medidas, por qué el perfil radial de atenuación resultó insuficiente para reconstruir la pérdida integrada en el anillo.

El residuo del 0.0548\,\% y la escasa variación de los términos durante el último período respaldaron el cierre local del balance para esta geometría y esta malla. La fracción disipada caracterizó la energía que atravesó el anillo seleccionado y dependió de sus límites. El residuo cuantificó la discrepancia de cierre entre integrales calculadas en la misma corrida. La repetición a 10 px/mm redujo la fracción disipada en 0.337943 puntos porcentuales. La diferencia relativa de 0.4181\,\% respecto de la malla fina fue pequeña frente al cambio cercano al 30\,\% de las potencias absolutas. Este contraste mostró que la normalización absorbió buena parte de la dependencia de la amplitud con la malla. La representación de un conductor de dos celdas de diámetro y de una fuente de dimensiones reducidas ofreció una explicación plausible de esa sensibilidad, cuya contribución individual quedó sin aislar.

\section{Conservación local y alcance de la validación física}
La disipación integrada conservó el orden observado en la atenuación de los medios con pérdidas. Las fracciones de 81.17\,\%, 74.53\,\% y 30.38\,\% en agua salada, músculo y concreto describieron el destino de la potencia que entró en cada anillo (Cuadro~\ref{tab:balance-medios}). Los límites se expresaron en longitudes de onda y correspondieron a volúmenes físicos diferentes. La comparación quedó, por ello, referida a regiones eléctricamente semejantes. La interpretación como eficiencia total de radiación o como fracción de la potencia suministrada por el generador quedó fuera del alcance de estas superficies de control.

En vacío, la conductividad nula impuso la ausencia de pérdida de Joule. La comprobación adicional consistió en contrastar de forma independiente los flujos entrantes y salientes. El residuo final de 0.013943\,\% mostró una concordancia estrecha entre ambos, aunque aumentó desde 0.000308\,\% durante el último intervalo. La máxima variación temporal de un término fue de 0.017448\,\% de la entrada final, comparable en magnitud con ese residuo. En consecuencia, la comprobación temporal exigió mayor cautela en vacío. La conservación local aportó un control físico complementario al ajuste radial, cuyo criterio de calidad había fallado.

En concreto, agua salada y músculo, los residuos finales de 0.002011\,\%, 0.054794\,\% y 0.032973\,\% se acompañaron de cambios temporales de los términos inferiores a 0.000144\,\% de la entrada. Esta separación de magnitudes respaldó la estabilidad de las integrales en los dos instantes examinados. En agua salada, el descenso del residuo al refinar la malla también resultó compatible con una contribución de la discretización y de la cuadratura. El almacenamiento transitorio de energía, la interpolación de campos y las integrales de superficie y volumen pudieron contribuir simultáneamente al cierre observado. La disponibilidad de dos muestras, tomadas a una separación cercana a un período, limitó la evaluación de transitorios más lentos.

\section{Incertidumbres y alcance de la resolución adoptada}
Se distinguieron tres niveles de evaluación. El remuestreo por bloques cuantificó la incertidumbre condicional del ajuste de atenuación. La comparación con la solución analítica del medio homogéneo midió la discrepancia física de la pendiente aparente. Las diferencias entre mallas describieron la sensibilidad del procedimiento numérico completo, incluidos el muestreo y la selección del modelo de ajuste. La ausencia de un orden asintótico estable dejó sin determinar una cota del error de discretización al continuo. Para el balance se añadieron los diagnósticos de cierre y de variación temporal, y quedó pendiente una estimación global de incertidumbre que integrara esas contribuciones.

La resolución común de 5 px/mm se sustentó en una precisión y un costo explícitos. Permitió recuperar el orden de atenuación de los medios con pérdidas, con discrepancias teóricas entre 1.31\,\% y 5.12\,\%, y diferencias frente a las mallas más finas disponibles de 0.26\,\% a 3.59\,\%. En el balance de agua salada, el refinamiento a 10 px/mm modificó la fracción disipada en 0.42\,\% respecto del valor fino. A la vez, la dependencia de las potencias absolutas y la representación limitada del conductor restringieron la precisión atribuible a la malla común.

El costo medido reforzó esa delimitación. El vacío consumió 4.75 GiB de memoria residente y 7.49 h a 5 px/mm, dentro de un entorno WSL con 7.44 GiB de memoria total (Cuadro~\ref{tab:balance-costo}). Un refinamiento sistemático de todos los medios habría requerido recursos adicionales. Así, la comparación a 5 px/mm se justificó para las tendencias y los cocientes documentados, con los márgenes observados en cada magnitud. La validación alcanzada combinó concordancia teórica de la atenuación en los medios con pérdidas, sensibilidad de malla y cierre energético local. Su extensión a potencias absolutas convergidas, a la eficiencia global de la antena y a otros parámetros materiales quedó pendiente de verificaciones específicas.

\section{Pérdidas de propagación y evaluación de los objetivos}
Los perfiles de las Figuras~\ref{fig:perfiles-vacio-concreto} y~\ref{fig:perfiles-agua-musculo} permitieron examinar el decaimiento y la estructura de los residuos en la malla común. Las seis figuras del Anexo~\ref{an:fdtd-perfiles} respaldaron las comparaciones de refinamiento. En agua salada se observó una aproximación progresiva a la pendiente teórica. En músculo y concreto se cuantificó la sensibilidad con las resoluciones disponibles. En vacío persistió la limitación del ajuste, aunque el balance aportó una comprobación independiente del transporte de potencia.

Para el modelo M0, la pendiente ajustada permitió expresar una pérdida radial relativa dentro de la ventana mediante
\begin{equation}
 L(r,r_0)=20\log_{10}(r/r_0)+\alpha_{\rm app}(r-r_0),\qquad r,r_0\text{ en mm}.
\label{eq:perdida-relativa-modelo}
\end{equation}
Esta expresión se obtuvo al deshacer la compensación por $r^2$. Para un modelo con corrección $g(r)$ se añadió $g(r_0)-g(r)$; en M1 esa función fue sinusoidal. La relación describió el cociente del flujo radial observado entre dos posiciones de la misma configuración. Su uso quedó limitado a la ventana del ajuste y a la calidad registrada.

El objetivo general y el primer objetivo específico se atendieron en el alcance de propagación dentro de medios homogéneos con la antena inmersa. Las pérdidas de penetración a través de una interfaz y la transmisión por espesores de una pared requirieron geometrías adicionales. La definición de regiones, unidades, normalización y controles de calidad constituyó una propuesta de procedimiento para el segundo objetivo específico.

Respecto del tercer objetivo, los resultados permitieron identificar un término de atenuación para una representación radial relativa. El ajuste sobre las mismas muestras evaluó concordancia interna. La capacidad predictiva frente a datos independientes y la calibración de modelos de enlace quedaron pendientes. El cambio de modelo entre mallas, la extracción deficiente en vacío y la dependencia de las potencias absolutas delimitaron esa aplicación.

El umbral interno de sensibilidad del 5\,\% se cumplió para las pendientes con refinamiento disponible y para la fracción disipada de agua salada. Las potencias absolutas de ese medio superaron el umbral al cambiar cerca del 30\,\%. Se justificó así el uso comparativo de las pendientes y del cociente energético dentro de los límites medidos. El residuo pequeño del balance describió la conservación local, y la evaluación de exactitud requirió considerar además la sensibilidad espacial y temporal.
